% !PRINTED 334
\weberpage{334}
Bleiben wir zunächst bei der ersten Differenz stehen, so ergiebt sich
als \gesperrt{erste Correction}
\weq{7}{ u = \delta\frac{\varDelta}{\varDelta_\alpha}, }
und wenn wir nun
\[ u = \delta\frac{\varDelta}{\varDelta_\alpha} + u' \]
setzen, so folgt aus (6), wenn man im zweiten Gliede $u'$ weglässt,
\[ 0 = \frac{\varDelta_\alpha}{\delta}u' +
  \frac{\varDelta'_\alpha \varDelta(\varDelta-\varDelta_\alpha)}{2\varDelta_\alpha^2}, \]
also als zweite Correction:
\weq{8}{ u' = \delta\frac{\varDelta'_\alpha
  \varDelta(\varDelta_\alpha-\varDelta)}{2\varDelta_\alpha^3}. }

Die erste Correction erhält man dadurch, dass man den zwischen $\alpha$
und $\alpha+\delta$ verlaufenden Curvenbogen durch die Sehne ersetzt,
wie oben, die zweite dadurch, dass man den Bogen zwischen $\alpha,
\alpha+\delta, \alpha+2\delta$ durch eine Parabel ersetzt, die durch
dieselben Punkte geht, die aber jetzt ihre \gesperrt{Axe mit der $x$-Axe parallel}
hat.

Nehmen wir als Beispiel die Gleichung
\[ f(x) = x^3 - 2x - 2 = 0, \]
die zwischen 1,7 und 1,8 eine reelle Wurzel hat.

Man berechnet
\[
\begin{array}{llll}
x = 1,7, & f(x) = -0,487, & \varDelta_\alpha = 0,719, & \varDelta'_\alpha = 0,108,\\
\phantom{x = {}}1,8, & \phantom{f(x) {}}{=} +0,232, & \phantom{\varDelta_\alpha {}}{=} 0,827, & \\
\phantom{x = {}}1,9, & & & \phantom{\varDelta'_\alpha {}}{=} 1,059.
\end{array}
\]
Da $f(x) = 0$ sein soll, so ist
\[ \varDelta = 0,487 \]
zu setzen und die erste Correction zu $\alpha = 1,7$ ist nach (7)
\[ u = 0,06773, \]
die zweite Correction ergiebt nach (8)
\[ u' = 0,0164, \]
also
\[ \alpha + u + u' = 1,76937. \]

Der auf andere Weise berechnete genauere Werth ist 1,76929\ldots\ Wir
haben also ein in den ersten drei Decimalen genaues Resultat. Wir haben
aber hier kein anderes Mittel, um die Genauig-
%CONT>>
